\ExamPart{Trắc nghiệm trả lời ngắn}{2}{Thí sinh trả lời từ câu 1 đến câu 4.}

\NQuestion Một xưởng may bán mỗi áo với giá $300$ nghìn đồng. Chi phí sản xuất $x$ chiếc áo là
\[
C(x)=x^2+80x+500 \quad (\text{nghìn đồng}).
\]
Hỏi xưởng phải bán ít nhất bao nhiêu chiếc áo để không bị lỗ?
\AnswerLine{$3$}
\SolutionBlock{Doanh thu là $R(x)=300x$ (nghìn đồng). Điều kiện không lỗ:
\[
300x\ge x^2+80x+500 \Leftrightarrow x^2-220x+500\le0.
\]
Phương trình bậc hai tương ứng có nghiệm
\[
x=110\pm10\sqrt{116}.
\]
Ta có $110-10\sqrt{116}\approx2{,}30$. Vì $x$ là số nguyên dương, số áo ít nhất phải bán là $3$.}

\NQuestion Trong mặt phẳng tọa độ $Oxy$, một điểm cấp nước ở vị trí $M(4;1)$ và bờ sông được mô hình hoá bởi đường thẳng $d:2x-y-3=0$. Tính khoảng cách ngắn nhất từ $M$ đến bờ sông.
\AnswerLine{$\dfrac{4}{\sqrt5}=\dfrac{4\sqrt5}{5}$}
\SolutionBlock{Khoảng cách từ $M$ đến $d$ là
\[
d(M,d)=\frac{|2\cdot4-1-3|}{\sqrt{2^2+(-1)^2}}=\frac{4}{\sqrt5}=\frac{4\sqrt5}{5}.
\]}

\NQuestion Phương trình $x^4-13x^2+36=0$ có bao nhiêu nghiệm thực phân biệt?
\AnswerLine{$4$}
\SolutionBlock{Đặt $t=x^2\ (t\ge0)$, ta được $t^2-13t+36=0 \Leftrightarrow (t-4)(t-9)=0 $. Suy ra $t=4$ hoặc $t=9$, nên $x=\pm2,\pm3$. Vậy có $4$ nghiệm thực phân biệt.}

\NQuestion (Câu khó) Có bao nhiêu số nguyên $m$ để phương trình $|x^2-2x-3|=m$ có đúng $3$ nghiệm thực phân biệt?
\AnswerLine{$1$}
\SolutionBlock{Ta có
\[
x^2-2x-3=(x-1)^2-4.
\]
Phương trình đã cho tương đương với
\[
(x-1)^2-4=m \quad \text{hoặc} \quad (x-1)^2-4=-m,
\]
hay
\[
(x-1)^2=m+4 \quad \text{hoặc} \quad (x-1)^2=4-m.
\]
Vì vế trái không âm nên cần $m\ge0$.
\begin{itemize}[leftmargin=1.5em]
\item Nếu $m=0$ thì hai phương trình trùng nhau: $(x-1)^2=4$, có $2$ nghiệm.
\item Nếu $0<m<4$ thì mỗi phương trình cho $2$ nghiệm phân biệt và không trùng nhau, tổng cộng $4$ nghiệm.
\item Nếu $m=4$ thì ta được $(x-1)^2=8$ hoặc $(x-1)^2=0$, tức có $2+1=3$ nghiệm phân biệt.
\item Nếu $m>4$ thì phương trình thứ hai vô nghiệm, phương trình thứ nhất có $2$ nghiệm.
\end{itemize}
Vậy chỉ có đúng một số nguyên $m$, đó là $m=4$.}
